\section{Cái trần, và cách chứng minh nó}
\label{sec:ch13-tran-va-chung-minh}

\index{giới hạn lý thuyết!khoảng cách độ phức tạp|(}%
Mục trước để lại một lời hứa: ô chữ nhật ở góc dưới bên trái của
hình~\ref{fig:ch13-mien-kha-thi} là rỗng với những mô hình đủ phức tạp. Mục này
chứng minh điều đó, và cái đáng chú ý không phải kết luận mà là chứng minh dài
bao nhiêu.

Trước hết, đầu mút bên phải của đường cong. Bài báo chứng minh rằng
$\kappa_f(0) = K(f)$: số bit ít nhất đủ để giải thích $f$ không sai chút nào
đúng bằng \tn{độ phức tạp Kolmogorov}{Kolmogorov complexity} của chính
$f$~\autocite{p26limits}. Chiều dễ là
chiều lên: bản thân $f$ đạt sai số không và tốn $K(f)$ bit. Chiều xuống là nội
dung thật: nếu một hàm $g$ với $K(g) < K(f)$ đạt sai số không thì $g$ trùng $f$
trên mọi đầu vào, nên chương trình tính $g$ cũng tính $f$, nên
$K(f) \le K(g) < K(f)$, vô lý. Nói bằng lời thì kết quả này phát biểu rằng không
có bản rút gọn nào giải thích trọn vẹn một mô hình; muốn không sai chỗ nào thì
phải chép lại cả mô hình.

Bỏ đòi hỏi \enquote{không sai chút nào} thì được kết quả trung tâm của bài.

\begin{theorem}[Khoảng cách độ phức tạp]
\label{thm:ch13-khoang-cach-do-phuc-tap}
Tồn tại một hằng số $c$ chỉ phụ thuộc \tn{máy Turing phổ dụng}{universal Turing
machine} đã chọn, sao cho với
mọi mô hình $f$ và mọi hàm $g$ thỏa $K(g) < K(f) - c$, tồn tại một đầu vào
$x \in \mathcal{X}$ với $f(x) \ne g(x)$.
\end{theorem}

\begin{proof}
Giả sử ngược lại rằng $f(x) = g(x)$ với mọi $x \in \mathcal{X}$, tức $f = g$.
Gọi $p_g$ là chương trình ngắn nhất tính $g$, độ dài $K(g)$. Vì $f = g$ nên
chính $p_g$ cũng tính $f$, do đó $K(f) \le |p_g| = K(g)$. Điều này mâu thuẫn với
giả thiết $K(g) < K(f) - c$.
\end{proof}

Định lý~\ref{thm:ch13-khoang-cach-do-phuc-tap} là định lý 2.23 của bài
báo~\autocite{p26limits}, và chứng minh trên là chứng minh của bài, viết đủ chứ
không rút gọn. Nó dài bốn dòng, không dùng một giả thiết nào về $f$, và không
mượn kết quả nào ngoài định nghĩa của $K$. Cả hai mặt của kết quả nằm ở đúng chỗ
đó. Vì không giả thiết gì nên nó đúng cho mọi mô hình, mọi kiến trúc, mọi bộ dữ
liệu và mọi phương pháp giải thích đã có hay sẽ có; không có tiến bộ kỹ thuật
nào lách được nó. Nhưng cũng vì không giả thiết gì nên nó nói rất ít: nó khẳng
định có một đầu vào mà hai hàm cho kết quả khác nhau, không nói có bao nhiêu đầu
vào như thế, chúng nằm ở đâu, hay chênh lệch lớn tới đâu. Một lời giải thích
lệch đúng một điểm, và lệch một lượng nhỏ tùy ý, vẫn thỏa định lý.

\index{không suy biến}%
Bịt lỗ hổng ấy cần thêm một điều kiện, và bài đặt tên cho nó. Một hàm $f$ được
gọi là \tn{không suy biến}{non-degenerate} ở mức $\delta$ nếu mọi đầu ra khả dĩ
khác $f(x)$ đều cách $f(x)$ một khoảng lớn hơn $\delta$, với mọi
$x$~\autocite{p26limits}. Điều kiện này nói rằng các đầu ra của $f$ tách nhau đủ
xa để một chỗ lệch là một chỗ nhìn thấy được: nếu $g(x) \ne f(x)$ thì tất yếu
$\rho(f(x),g(x)) > \delta$.

Vì sao phải phát biểu riêng điều đó thì thấy rõ qua hai trường hợp. Với một bộ
phân loại có nhãn rời rạc và khoảng cách giữa hai nhãn bất kỳ là một hằng số,
điều kiện tự động đúng với mọi $\delta$ nhỏ hơn hằng số ấy. Với một mô hình hồi
quy nhận giá trị trên toàn trục số thì nó hỏng: hai đầu ra có thể cách nhau
$10^{-9}$, nên một hàm $g$ lệch khỏi $f$ ở rất nhiều điểm vẫn giữ được sai số
tổng cộng nhỏ hơn $\delta$. Bài gọi những chỗ lệch ấy là các bất đồng vô hình,
và mục~\ref{sec:ch13-tran-noi-gi} sẽ dùng lại chúng.

Cộng hai điều lại thì được phát biểu mà chương này thật sự cần. Nếu $f$ không
suy biến ở mức $\delta$ và $K(g) < K(f) - c$ thì
$\mathcal{E}_{\infty}(f,g) > \delta$~\autocite{p26limits}. Đường cong trong
hình~\ref{fig:ch13-mien-kha-thi} vì thế không chỉ dương mà còn nằm trên ngưỡng
$\delta$ ở mọi ngân sách nhỏ hơn $K(f) - c$, và ô chữ nhật đúng là rỗng. Ở đầu
kia của thang, bài ghi trường hợp tầm thường cho đủ bộ: một hàm với
$K(f) \le C$ thì giải thích được hoàn hảo trong đúng $C$ bit, vì chính nó là một
lời giải thích hợp lệ~\autocite{p26limits}.

\index{bộ ba bất khả|(}%
Bài báo dành mười lăm trang trong sáu mươi lăm, mục 4 cùng toàn bộ phụ lục, cho một chủ
đề mà tên bài không báo trước: quy định pháp lý về \tn{khả năng giải
thích}{explainability}. Bài đọc
ba khung cụ thể theo hướng ấy, gồm điều 13 của EU AI Act, các yêu cầu về mô hình
trong ngành tài chính và hướng dẫn của FDA cho thiết bị y tế dùng học máy, và
lập luận rằng mỗi khung đều đòi cùng lúc lời giải thích đọc được và độ khớp cao
mà không thừa nhận sự đánh đổi giữa hai thứ~\autocite{p26limits}. Sách không đi
vào phần ấy, vì đối tượng của cuốn sách này là dụng cụ đo của lĩnh vực chứ không
phải chính sách. Một kết quả trong đó thì phải nêu, vì nó là định lý vừa chứng
minh viết lại bằng thứ ngôn ngữ khác.

Bài gọi cấu trúc lôgic ấy là \tn{bộ ba bất khả}{trilemma}, cùng dạng với bộ ba
bất khả trong kinh tế học quốc tế và với định lý CAP của hệ phân
tán~\autocite{p26limits}.

\begin{theorem}[Bộ ba bất khả trong quy định]
\label{thm:ch13-bo-ba-bat-kha}
Xét một khung quy định áp cùng lúc ba đòi hỏi lên các hệ thống không suy biến ở
mức $\delta$. Đòi hỏi thứ nhất cho phép triển khai hệ thống $f$ với $K(f)$ lớn
tùy ý. Đòi hỏi thứ hai buộc mỗi hệ thống được triển khai phải kèm một lời giải
thích $g$ với $K(g) \le b_{\text{người}}$, trong đó $b_{\text{người}}$ là một
hằng số cố định biểu thị sức đọc của con người. Đòi hỏi thứ ba buộc sai số của
lời giải thích thỏa $\mathcal{E}(f,g) \le \delta$. Khi ấy hai đòi hỏi bất kỳ
trong ba đòi hỏi trên thỏa được đồng thời, cả ba thì không, và không đòi hỏi nào
suy ra được từ hai đòi hỏi còn lại.
\end{theorem}

Đây là định lý 4.6 của bài~\autocite{p26limits}. Vế \enquote{cả ba thì không}
được chứng minh trong phụ lục, và đọc phần chứng minh ấy thì thấy hai bước cốt
lõi của nó là định lý~\ref{thm:ch13-khoang-cach-do-phuc-tap} cộng với điều kiện
không suy biến, thay $K(g)$ bằng $b_{\text{người}}$. Đòi hỏi thứ nhất cho phép
chọn một $f$ với $K(f) \ge b_{\text{người}} + c$; đòi hỏi thứ hai ép
$K(f) - K(g) \ge c$; định lý cho một điểm lệch; và điều kiện không suy biến biến
điểm lệch ấy thành một sai số lớn hơn $\delta$, mâu thuẫn với đòi hỏi thứ ba.
Vế \enquote{hai đòi hỏi bất kỳ thì được} chứng minh bằng cách dựng ví dụ cho
từng cặp: bỏ đòi hỏi thứ ba thì lấy $g$ là hàm hằng, bỏ đòi hỏi thứ hai thì lấy
$g$ bằng chính $f$, bỏ đòi hỏi thứ nhất thì chỉ cho phép những mô hình dưới
$b_{\text{người}}$ bit.

Nhận xét của sách về kết quả này, không phải của bài: định
lý~\ref{thm:ch13-bo-ba-bat-kha} không thêm toán học nào vào định
lý~\ref{thm:ch13-khoang-cach-do-phuc-tap}. Cái nó thêm là một cách đọc, và cách
đọc ấy có giá trị riêng, vì nó chỉ ra rằng ba thứ mà một bên soạn quy định
thường muốn cùng lúc là ba thứ đã được chứng minh không đi cùng nhau được. Bài
tự nhấn mạnh rằng đây là một kết quả bất khả có điều kiện: nó không nói rằng AI
phức tạp là cần thiết, cũng không kê đơn nên bỏ đòi hỏi nào, mà chỉ làm rõ hệ
quả toán học của từng lựa chọn~\autocite{p26limits}.

\index{bộ ba bất khả|)}%
\index{giới hạn lý thuyết!khoảng cách độ phức tạp|)}%
